\chapter{Peluang dan Penyampelan}\label{ch:g10:proba}

\omterm{def:g10:proba:distribution}{Peluang} memberikan bilangan kepada kebetulan. Berangkat dari percobaan dengan
hasil yang sama-sama mungkin --- dadu, koin, kartu --- bab ini membangun
aturan dasar untuk menghitung \omterm{def:g10:proba:distribution}{peluang} sebuah \omterm{def:g10:proba:events}{kejadian}, memperkenalkan \omterm{met:g10:proba:tree}{diagram
pohon} untuk percobaan dua langkah, lalu berakhir pada kaitan antara
\omterm{def:g10:proba:distribution}{peluang} dan \omterm{def:g10:stats:series}{frekuensi relatif} yang teramati. Teorinya berlanjut pada
\cref{ch:g11:prob}.

\section{Percobaan acak dan kejadian}

\begin{definition}[Hasil dan kejadian]\label{def:g10:proba:events}
Sebuah \emph{percobaan acak}\index{percobaan acak} adalah percobaan yang
hasilnya tidak dapat diramalkan (melempar dadu). Himpunan semua hasil yang
mungkin disebut \emph{ruang sampel}\index{ruang sampel}
$\Omega = \{\omega_1, \dots, \omega_n\}$; anggotanya disebut
\emph{hasil}. Sebuah \emph{kejadian}\index{kejadian} adalah himpunan hasil;
sebuah kejadian \emph{terjadi} apabila hasil percobaannya termasuk di dalamnya.
\end{definition}

\begin{definition}[Distribusi peluang]\label{def:g10:proba:distribution}
Memberikan kepada setiap hasil $\omega_i$ sebuah bilangan $p_i \geq 0$, dengan
$p_1 + \dots + p_n = 1$, mendefinisikan sebuah \emph{peluang}\index{peluang}:
peluang $\P(A)$ sebuah \omterm{def:g10:proba:events}{kejadian} $A$ adalah jumlah $p_i$ dari
hasil-hasil di dalam $A$. Khususnya $\P(\Omega) = 1$ dan \omterm{def:g10:proba:events}{kejadian} mustahil
$\varnothing$ berpeluang $0$.
\end{definition}

\begin{proposition}[Hasil yang sama-sama mungkin]\label{prop:g10:proba:uniform}
Jika $n$ hasilnya sama-sama mungkin (masing-masing $p_i = \frac1n$), maka untuk
setiap \omterm{def:g10:proba:events}{kejadian} $A$:
\[
\P(A) = \frac{\text{banyaknya hasil dalam } A}{n}
= \frac{\text{kasus yang menguntungkan}}{\text{kasus yang mungkin}} .
\]
\end{proposition}

\begin{proof}
$A$ memuat sejumlah $k$ hasil, masing-masing berpeluang $\frac1n$,
sehingga $\P(A) = \frac kn$.
\end{proof}

\begin{example}\label{ex:g10:proba:die}
Lemparlah sebuah dadu setimbang: $\Omega = \{1, 2, 3, 4, 5, 6\}$, setiap
hasilnya berpeluang $\frac16$. \omterm{def:g10:proba:events}{Kejadian} $A$: ``hasilnya genap'' adalah
$A = \{2, 4, 6\}$, jadi $\P(A) = \frac36 = \frac12$. \omterm{def:g10:proba:events}{Kejadian} $B$: ``hasilnya
paling sedikit $5$'' adalah $B = \{5, 6\}$, jadi $\P(B) = \frac26 =
\frac13$.
\end{example}

\section{Memadukan kejadian}

\begin{definition}[Irisan, gabungan, komplemen]\label{def:g10:proba:operations}
Misalkan $A$ dan $B$ dua \omterm{def:g10:proba:events}{kejadian}.
\begin{itemize}
  \item $A \cap B$ (``$A$ dan $B$'') terjadi apabila keduanya terjadi;
  \item $A \cup B$ (``$A$ atau $B$'') terjadi apabila paling sedikit satu terjadi;
  \item \emph{komplemen}\index{komplemen} $\bar A$ (``bukan $A$'')
        terjadi tepat ketika $A$ tidak terjadi;
  \item $A$ dan $B$ disebut \emph{saling lepas}\index{saling lepas} apabila keduanya tidak dapat terjadi
        bersamaan: $A \cap B = \varnothing$.
\end{itemize}
\end{definition}

\begin{theorem}[Aturan penjumlahan]\label{thm:g10:proba:addition}
Untuk semua \omterm{def:g10:proba:events}{kejadian} $A$ dan $B$:
\[
\P(A \cup B) = \P(A) + \P(B) - \P(A \cap B),
\qquad
\P(\bar A) = 1 - \P(A) .
\]
Untuk \omterm{def:g10:proba:events}{kejadian} yang \omterm{def:g10:proba:operations}{saling lepas}, aturan pertamanya menyederhana menjadi
$\P(A \cup B) = \P(A) + \P(B)$.
\end{theorem}

\begin{proof}
Menjumlahkan $\P(A) + \P(B)$ menambahkan \omterm{def:g10:proba:distribution}{peluang} setiap hasil dalam
$A \cup B$ satu kali, \emph{kecuali} hasil dalam $A \cap B$, yang terhitung
dua kali --- sekali di $A$ dan sekali di $B$. Mengurangkan $\P(A \cap B)$
memperbaiki penghitungan ganda itu. Untuk \omterm{def:g10:proba:operations}{komplemennya}: $A$ dan $\bar A$ \omterm{def:g10:proba:operations}{saling
lepas} dan $A \cup \bar A = \Omega$, jadi $\P(A) + \P(\bar A) = \P(\Omega) = 1$.
\end{proof}

\begin{omfigure}
\begin{tikzpicture}[scale=0.9]
  \draw[omExo, thick, rounded corners] (-3.1,-1.7) rectangle (3.4,1.7);
  \node[omExo, font=\small] at (3.0,1.45) {$\Omega$};
  \begin{scope}
    \clip (-0.8,0) circle (1.25);
    \fill[omProp!25] (0.8,0) circle (1.25);
  \end{scope}
  \draw[omDef, thick] (-0.8,0) circle (1.25);
  \draw[omThm, thick] (0.8,0) circle (1.25);
  \node[omDef, font=\small] at (-1.6,0.9) {$A$};
  \node[omThm, font=\small] at (1.7,0.9) {$B$};
  \node[font=\small] at (0,0) {$A \cap B$};
\end{tikzpicture}

{\small Dua \omterm{def:g10:proba:events}{kejadian} yang bertumpang tindih: pada $\P(A) + \P(B)$, \omterm{def:g10:numbers:interunion}{irisan}
$A \cap B$ yang terarsir terhitung dua kali, dan itulah yang menjelaskan
pengurangan pada aturan penjumlahannya.}
\end{omfigure}

\begin{example}\label{ex:g10:proba:cards}
Ambil satu kartu dari satu set $52$ kartu baku. Misalkan $A$: ``kartunya
bergambar hati'' ($13$ kartu) dan $B$: ``kartunya raja'' ($4$ kartu). Maka
$A \cap B$ adalah ``raja hati'' ($1$ kartu), dan
\[
\P(A \cup B) = \frac{13}{52} + \frac{4}{52} - \frac{1}{52} = \frac{16}{52}
= \frac{4}{13}.
\]
Aturan \omterm{def:g10:proba:operations}{komplemen} sering menjadi jalan tercepat: \omterm{def:g10:proba:distribution}{peluang} bahwa kartunya
\emph{bukan} hati adalah $1 - \frac{13}{52} = \frac34$.
\end{example}

\section{Diagram pohon}

\begin{method}[Percobaan dua langkah]\label{met:g10:proba:tree}
Untuk percobaan yang dikerjakan dalam dua langkah:
\begin{enumerate}
  \item gambarlah sebuah \emph{pohon}\index{diagram pohon}: satu cabang untuk
        setiap hasil yang mungkin pada langkah pertama, lalu dari
        masing-masingnya, satu cabang untuk setiap hasil langkah kedua, sambil
        menuliskan \omterm{def:g10:proba:distribution}{peluang} pada setiap cabang (cabang yang meninggalkan sebuah
        simpul harus berjumlah $1$);
  \item \omterm{def:g10:proba:distribution}{peluang} sebuah \emph{jalur} (sehelai daun \omterm{met:g10:proba:tree}{pohon} itu) adalah
        hasil kali \omterm{def:g10:proba:distribution}{peluang} sepanjang cabangnya;
  \item \omterm{def:g10:proba:distribution}{peluang} sebuah \omterm{def:g10:proba:events}{kejadian} adalah jumlah \omterm{def:g10:proba:distribution}{peluang}
        jalur yang mewujudkannya.
\end{enumerate}
\end{method}

\begin{example}\label{ex:g10:proba:tree}
Sebuah guci berisi $3$ bola merah dan $2$ bola biru. Ambil satu bola,
kembalikan, lalu ambil lagi. Setiap pengambilan memberi merah (M) berpeluang
$\frac35$ dan biru (B) berpeluang $\frac25$.

\begin{omfigure}
\begin{tikzpicture}[scale=1.0]
  \coordinate (root) at (0,0);
  \coordinate (R) at (2.3,1.1);
  \coordinate (B) at (2.3,-1.1);
  \coordinate (RR) at (5.0,1.7);
  \coordinate (RB) at (5.0,0.5);
  \coordinate (BR) at (5.0,-0.5);
  \coordinate (BB) at (5.0,-1.7);
  \draw (root) -- (R) node[midway, above, font=\small, sloped] {$\frac35$};
  \draw (root) -- (B) node[midway, below, font=\small, sloped] {$\frac25$};
  \draw (R) -- (RR) node[midway, above, font=\small, sloped] {$\frac35$};
  \draw (R) -- (RB) node[midway, below, font=\small, sloped] {$\frac25$};
  \draw (B) -- (BR) node[midway, above, font=\small, sloped] {$\frac35$};
  \draw (B) -- (BB) node[midway, below, font=\small, sloped] {$\frac25$};
  \node[omThm, font=\small] at (R) [above left=-2pt] {M};
  \node[omDef, font=\small] at (B) [below left=-2pt] {B};
  \node[right, font=\small] at (RR)
    {MM:\ $\frac35 \times \frac35 = \frac{9}{25}$};
  \node[right, font=\small] at (RB)
    {MB:\ $\frac35 \times \frac25 = \frac{6}{25}$};
  \node[right, font=\small] at (BR)
    {BM:\ $\frac25 \times \frac35 = \frac{6}{25}$};
  \node[right, font=\small] at (BB)
    {BB:\ $\frac25 \times \frac25 = \frac{4}{25}$};
\end{tikzpicture}

{\small \omterm{met:g10:proba:tree}{Pohon} dua pengambilan dengan pengembalian. Keempat \omterm{def:g10:proba:distribution}{peluang} jalurnya
berjumlah $1$.}
\end{omfigure}

\omterm{def:g10:proba:distribution}{Peluang} memperoleh dua bola berwarna sama adalah
$\P(\text{MM}) + \P(\text{BB}) = \frac{9}{25} + \frac{4}{25} =
\frac{13}{25}$; \omterm{def:g10:proba:distribution}{peluang} memperoleh paling sedikit satu biru adalah
$1 - \P(\text{MM}) = 1 - \frac{9}{25} = \frac{16}{25}$ (aturan \omterm{def:g10:proba:operations}{komplemen}).
\end{example}

\begin{remark}
Tanpa pengembalian, \omterm{def:g10:proba:distribution}{peluang} pengambilan kedua bergantung pada yang
pertama: dengan guci di atas, setelah mengambil sebuah bola merah, hanya
tersisa $2$ merah dan $2$ biru, sehingga cabang keduanya membawa $\frac24$ dan
$\frac24$. Metode \omterm{met:g10:proba:tree}{pohonnya} bekerja tanpa perubahan; hanya bilangan pada cabang
tingkat keduanya yang berubah.
\end{remark}

\section{Frekuensi relatif dan penyampelan}

Ketika sebuah \omterm{def:g10:proba:events}{percobaan acak} diulang berkali-kali, \omterm{def:g10:stats:series}{frekuensi relatif} yang
teramati untuk sebuah \omterm{def:g10:proba:events}{kejadian} mendekati \omterm{def:g10:proba:distribution}{peluangnya} --- itulah sebabnya
\omterm{def:g10:proba:distribution}{peluang} berguna untuk memerikan dunia nyata.

\begin{definition}[Sampel]\label{def:g10:proba:sample}
Sebuah \emph{sampel}\index{sampel} berukuran $n$ adalah hasil mengulang
percobaan yang sama sebanyak $n$ kali secara saling bebas.
\emph{\omterm{def:g10:stats:series}{Frekuensi relatif} teramati} sebuah \omterm{def:g10:proba:events}{kejadian} pada sampel itu adalah
banyaknya kemunculan dibagi $n$.
\end{definition}

\begin{example}\label{ex:g10:proba:sampling}
Sebuah koin setimbang memberi gambar dengan \omterm{def:g10:proba:distribution}{peluang} $0.5$. Melemparnya $100$
kali jarang memberi tepat $50$ gambar: \omterm{def:g10:proba:sample}{sampel} itu \emph{berfluktuasi}. Rangkaian
khas $100$ lemparan memberi \omterm{def:g10:stats:series}{frekuensi relatif} seperti $0.46$, $0.53$, $0.49$
--- dekat dengan $0.5$, tetapi tidak sama. \omterm{def:g10:proba:sample}{Sampel} yang lebih besar berfluktuasi
lebih kecil: dengan $10\,000$ lemparan, \omterm{def:g10:stats:series}{frekuensi relatif} yang teramati biasanya
berada dalam sekitar $0.01$ dari $0.5$.
\end{example}

\begin{remark}[Interval fluktuasi]
Sebuah aturan praktis yang berguna, yang dibenarkan pada \cref{ch:g11:binom}:
untuk \omterm{def:g10:proba:sample}{sampel} berukuran $n$ dari \omterm{def:g10:proba:events}{kejadian} berpeluang $p$ (dengan $n$ cukup
besar), \omterm{def:g10:stats:series}{frekuensi relatif} yang teramati terletak dalam \omterm{def:g10:numbers:interval}{interval}
\[
\intcc{p - \frac{1}{\sqrt n}}{\,p + \frac{1}{\sqrt n}}
\]
pada kira-kira $95\%$ \omterm{def:g10:proba:sample}{sampelnya}. Jika sebuah \omterm{def:g10:stats:series}{frekuensi relatif} teramati jatuh
jauh di luar \omterm{def:g10:numbers:interval}{interval} ini, ragukanlah nilai $p$ itu: inilah titik tolak
pengujian statistik.
\end{remark}

\begin{example}\label{ex:g10:proba:test}
Sebuah pabrik mengaku bahwa hanya $10\%$ suku cadangnya yang cacat
($p = 0.1$). Pada \omterm{def:g10:proba:sample}{sampel} berisi $n = 400$ suku cadang, $64$ di antaranya cacat:
\omterm{def:g10:stats:series}{frekuensi relatifnya} $\frac{64}{400} = 0.16$. \omterm{def:g10:numbers:interval}{Interval} fluktuasinya adalah
$\intcc{0.1 - \frac{1}{20}}{0.1 + \frac{1}{20}} = \intcc{0.05}{0.15}$,
dan $0.16$ jatuh di luarnya: \omterm{def:g10:proba:sample}{sampel} itu meragukan pengakuannya secara serius.
\end{example}

\section{Latihan}

\begin{exercise}[$\star$]\label{exo:g10:proba:1}
Sebuah dadu setimbang dilempar. Hitunglah \omterm{def:g10:proba:distribution}{peluang} setiap \omterm{def:g10:proba:events}{kejadian} berikut:
``memperoleh $6$''; ``memperoleh bilangan ganjil''; ``memperoleh paling banyak
$4$''; ``memperoleh $7$''.
\end{exercise}

\begin{exercise}[$\star$]\label{exo:g10:proba:2}
Sebuah kantong berisi $5$ kelereng hijau, $3$ kuning dan $2$ hitam; satu
kelereng diambil secara acak. Hitunglah \omterm{def:g10:proba:distribution}{peluang} bahwa kelereng itu hijau; bahwa
kelereng itu bukan hitam; bahwa kelereng itu hijau atau kuning.
\end{exercise}

\begin{exercise}[$\star$]\label{exo:g10:proba:3}
Di sebuah kelas berisi $30$ siswa, $18$ belajar bahasa Spanyol, $10$ belajar
bahasa Jerman, dan $4$ belajar keduanya. Seorang siswa dipilih secara acak.
Dengan aturan penjumlahan, hitunglah \omterm{def:g10:proba:distribution}{peluang} bahwa siswa itu belajar paling
sedikit satu di antara kedua bahasa itu, lalu \omterm{def:g10:proba:distribution}{peluang} bahwa ia tidak belajar keduanya.
\end{exercise}

\begin{exercise}[$\star$]\label{exo:g10:proba:4}
Sebuah koin setimbang dilempar dua kali. Gambarlah \omterm{met:g10:proba:tree}{pohon} percobaannya lalu
hitunglah \omterm{def:g10:proba:distribution}{peluang} memperoleh dua gambar; tepat satu gambar; paling sedikit satu
gambar.
\end{exercise}

\begin{exercise}[$\star\star$]\label{exo:g10:proba:5}
Sebuah guci berisi $4$ bola merah dan $6$ bola putih. Dua bola diambil satu
demi satu \emph{tanpa} pengembalian. Gambarlah \omterm{met:g10:proba:tree}{pohonnya} dengan
\omterm{def:g10:proba:distribution}{peluang} yang benar pada tingkat kedua, lalu hitunglah \omterm{def:g10:proba:distribution}{peluang} mengambil
dua bola merah, lalu \omterm{def:g10:proba:distribution}{peluang} mengambil dua bola berbeda warna.
\end{exercise}

\begin{exercise}[$\star\star$]\label{exo:g10:proba:6}
Dua dadu setimbang dilempar dan jumlah matanya dicatat.
\begin{enumerate}
  \item Jelaskan mengapa \omterm{def:g10:proba:events}{ruang sampelnya} dapat diambil sebagai $36$ pasangan
        yang sama-sama mungkin $(1,1), (1,2), \dots, (6,6)$.
  \item Hitunglah \omterm{def:g10:proba:distribution}{peluang} bahwa jumlahnya $7$; bahwa jumlahnya $12$;
        bahwa jumlahnya paling banyak $4$.
  \item Jumlah manakah yang paling mungkin?
\end{enumerate}
\end{exercise}

\begin{exercise}[$\star\star$]\label{exo:g10:proba:7}
Sebuah soal pilihan ganda menawarkan $4$ jawaban, satu di antaranya benar.
Seorang siswa menjawab $2$ soal semacam itu yang saling bebas, sepenuhnya
secara acak. Hitunglah \omterm{def:g10:proba:distribution}{peluang} menjawab benar keduanya, benar tepat satu, dan
benar tidak satu pun. Periksa bahwa ketiga \omterm{def:g10:proba:distribution}{peluang} itu berjumlah $1$.
\end{exercise}

\begin{exercise}[$\star\star$]\label{exo:g10:proba:8}
Seorang politikus mengaku memperoleh dukungan $50\%$. Dalam jajak pendapat
terhadap $n = 900$ pemilih yang dipilih acak, $396$ mendukungnya.
\begin{enumerate}
  \item Hitunglah \omterm{def:g10:stats:series}{frekuensi relatif} teramati dan \omterm{def:g10:numbers:interval}{interval} fluktuasinya
        di sekitar $p = 0.5$ untuk $n = 900$.
  \item Apakah jajak pendapat itu meragukan pengakuannya?
\end{enumerate}
\end{exercise}

\begin{exercise}[$\star\star\star$]\label{exo:g10:proba:9}
Sebuah permainan berbiaya $2$ untuk dimainkan: kamu melempar dadu setimbang
lalu menang $6$ bila muncul $6$, menang $3$ bila muncul $5$, dan tidak
memperoleh apa-apa selain itu.
\begin{enumerate}
  \item Hitunglah \omterm{def:g10:proba:distribution}{peluang} setiap jumlah kemenangan ($6$, $3$, $0$).
  \item Selama $600$ permainan, kira-kira berapa kali tiap hasil kamu harapkan?
        Hitunglah total harapan untung dan ruginya: apakah permainan ini layak
        dimainkan?
\end{enumerate}
\end{exercise}

\section{Soal: Tiga pintu yang memperdaya seribu
matematikawan}

\begin{problem}[{Soal akhir pekan --- paradoks Monty Hall dan
kerabatnya: bagaimana keterangan mengubah peluang, dan bagaimana sampel
memisahkan pengakuan dari angan-angan}]
\label{pb:g10:proba:1}
Pada tahun 1990 pengasuh rubrik Marilyn vos Savant menerbitkan sebuah teka-teki
tentang pemandu acara kuis, tiga pintu, sebuah mobil dan dua ekor kambing ---
lalu menjawabnya dengan benar. Sepuluh ribu pembaca, di antaranya ratusan
doktor matematika, menyurati redaksi untuk bersikeras bahwa ia keliru. Soal
ini mempersenjatai kamu dengan \omterm{met:g10:proba:tree}{pohon} (\cref{met:g10:proba:tree}),
\omterm{def:g10:proba:operations}{komplemen} dan aturan penjumlahan
(\cref{thm:g10:proba:addition}), membiarkanmu menuntaskan perdebatan itu
sendiri, lalu ditutup dengan alat yang menengahi setiap sengketa semacam itu
dalam praktik: \omterm{def:g10:proba:sample}{sampel} yang berfluktuasi.

\medskip
\noindent\textbf{Bagian I --- Aturannya, sebagai pemanasan.}
\begin{enumerate}
  \item Satu kartu dari satu set $52$ kartu baku: hitunglah
        $P(\text{raja atau hati})$ dengan aturan penjumlahan ---
        lalu katakan mengapa $\frac{4}{52} + \frac{13}{52}$ saja
        keliru.
  \item Tiga lemparan koin setimbang: hitunglah $P(\text{paling sedikit satu
        gambar})$ lewat \omterm{def:g10:proba:operations}{komplemennya}.
  \item Sebuah kantong berisi $3$ kelereng merah dan $2$ biru; dua diambil
        tanpa pengembalian. Dengan sebuah \omterm{met:g10:proba:tree}{pohon}, hitunglah
        $P(\text{dua merah})$ dan $P(\text{satu dari tiap warna})$.
  \item Lengkapkan distribusi pada pertanyaan 3 dengan
        $P(\text{dua biru})$ lalu periksa bahwa ketiga \omterm{def:g10:proba:distribution}{peluang} itu
        berjumlah $1$.
  \item Seseorang melaporkan, untuk dua \omterm{def:g10:proba:events}{kejadian} yang \emph{\omterm{def:g10:proba:operations}{saling lepas}},
        $P(A) = 0.6$ dan $P(B) = 0.5$. Buktikan kesalahannya dengan satu
        baris perhitungan.
\end{enumerate}

\medskip
\noindent\textbf{Bagian II --- Ketiga pintu itu.}
Permainannya: sebuah mobil bersembunyi di balik salah satu dari tiga pintu,
kambing di balik dua pintu lainnya. Kamu memilih satu pintu (misalnya pintu 1).
Pemandu acara --- yang \emph{tahu} di mana mobilnya --- membuka salah satu dari
dua pintu yang lain, selalu memperlihatkan kambing, lalu menawarimu kesempatan
untuk berpindah ke pintu tertutup yang tersisa.
\begin{enumerate}[resume]
  \item Sebelum menghitung apa pun: apakah berpindah itu menolong, merugikan,
        atau tidak berpengaruh? Tuliskan jawaban nalurimu --- dengan jujur.
  \item Sekarang \omterm{met:g10:proba:tree}{pohonnya}, menurut kasus letak mobil yang sebenarnya
        (masing-masing $\frac13$): jika mobilnya di balik pintumu,
        berpindah membuatmu kalah; jika mobilnya di balik salah satu pintu
        lain, apa yang harus dilakukan pemandu acara, dan apa yang lalu
        dimenangkan dengan berpindah?
        Hitunglah $P(\text{menang dengan berpindah})$.
  \item Versi satu barisnya: berpindah menang tepat ketika pilihan
        pertamamu keliru. Simpulkan sekali lagi.
  \item Kini ragam yang penting --- ``Monty Fall'': pemandu acaranya
        \emph{tidak} tahu di mana mobilnya, membuka salah satu
        pintu lain secara acak, dan karena beruntung memperlihatkan kambing.
        Daftarkan keenam skenario yang sama-sama mungkin (letak mobil $\times$
        pintu yang dibuka), buanglah skenario yang memperlihatkan mobilnya,
        lalu hitunglah \omterm{def:g10:proba:distribution}{peluang} bahwa berpindah itu menang
        \emph{di antara skenario yang tersisa}. Apa yang berubah, dan
        mengapa pengetahuan pemandu acara itu penting?
  \item Pemompa gerak hati: seratus pintu, kamu memilih satu, lalu
        pemandu acara yang tahu membuka $98$ pintu berkambing. Berpindah?
        Dengan \omterm{def:g10:proba:distribution}{peluang} menang berapa?
  \item Kamu berhadapan dengan seorang peragu yang, seperti seribu penulis
        surat itu, bersikeras ``50--50''. Perikan sebuah simulasi dengan kartu
        (tiga kartu, banyak ronde, seorang teman sebagai pemandu acara) lalu
        katakan berapa bagian kemenangan berpindah yang kamu harapkan setelah
        $300$ ronde --- dengan \omterm{def:g10:numbers:interval}{interval} fluktuasi bab ini
        sebagai batas kesalahanmu.
\end{enumerate}

\medskip
\noindent\textbf{Bagian III --- Kerabat paradoks itu.}
\begin{enumerate}[resume]
  \item Sebuah keluarga mempunyai dua anak; kamu mengetahui bahwa
        \emph{paling sedikit satu di antaranya laki-laki}. Daftarkan susunan
        dua anak yang sama-sama mungkin dan sesuai dengan keterangan itu, lalu
        hitunglah $P(\text{dua anak laki-laki})$.
  \item Keluarga yang sama, tetapi keterangannya kini: \emph{anak yang
        lebih tua laki-laki}. Hitunglah ulang. Mengapa kedua jawabannya
        berbeda, dan apa kesamaannya dengan kedua Monty tadi?
  \item Tiga koin setimbang: hitunglah $P(\text{ketiganya sama})$.
        Matematikawan d'Alembert berargumen pada tahun 1754 untuk jawaban
        yang lain: ``entah ketiganya sama entah tidak ---
        satu kasus dari dua.'' Sebutkan kekeliruannya secara tepat, dengan
        \cref{prop:g10:proba:uniform} di tangan.
  \item Sebuah permainan berbiaya $2$ euro: lemparlah dadu, menang $6$ euro
        bila muncul enam, $3$ euro bila muncul lima, dan tidak memperoleh
        apa-apa selain itu (\cref{exo:g10:proba:9}). Selama $600$ permainan,
        rata-rata untung per permainan berapa yang diramalkan penghitungan itu?
        Seorang penjamin asuransi memainkan matematika yang sama: \omterm{def:g10:proba:distribution}{peluang}
        $1\,\%$ per tahun untuk klaim sebesar $10\,000$ euro --- berapa
        premi yang \emph{adil}, dan mengapa premi yang sesungguhnya
        melampauinya?
\end{enumerate}

\medskip
\noindent\textbf{Bagian IV --- \omterm{def:g10:proba:sample}{Sampel} sebagai penengah.}
\begin{enumerate}[resume]
  \item Sebuah koin dilempar $100$ kali. Dengan memakai \omterm{def:g10:numbers:interval}{interval}
        fluktuasi $p \pm \frac{1}{\sqrt n}$, putuskan hasil mana di antara
        berikut yang patut mengangkat alis: $55$ gambar; $70$
        gambar.
  \item Sebuah jajak pendapat terhadap $1\,000$ pemilih memberi seorang
        calon $52\,\%$. Hitunglah \omterm{def:g10:numbers:interval}{interval} fluktuasi di sekitar $50\,\%$ lalu
        katakan apakah jajak pendapat itu \emph{membuktikan} bahwa sang calon
        unggul.
  \item Sebuah pabrik mengaku bahwa $4\,\%$ suku cadangnya cacat. \omterm{def:g10:proba:sample}{Sampel}
        berisi $400$ menunjukkan $28$ yang cacat ($7\,\%$); \omterm{def:g10:proba:sample}{sampel} berikutnya
        berisi $2\,500$ menunjukkan $175$ (lagi-lagi $7\,\%$). Ujilah
        pengakuannya terhadap masing-masing \omterm{def:g10:proba:sample}{sampel}. Mengapa putusannya berbeda?
  \item Dua sumber galat menghantui setiap survei: \emph{bias}
        penyampelan (soal akhir pekan tentang cara data berdusta di jilid
        sebelumnya) dan \emph{fluktuasi} penyampelan (bab ini). Mana yang
        mengecil ketika $n$ membesar, dan mana yang bertahan pada ukuran
        \omterm{def:g10:proba:sample}{sampel} berapa pun? Masing-masing satu kalimat.
  \item Penutup --- daftar periksa sang ahli \omterm{def:g10:proba:distribution}{peluang}, satu baris untuk
        setiap butir: hasil yang sama-sama mungkin sebelum dihitung
        (pertanyaan 15); \omterm{met:g10:proba:tree}{pohon} dan \omterm{def:g10:proba:operations}{komplemen} untuk percobaan bertahap dan
        untuk ``paling sedikit'' (pertanyaan 2--3); \emph{keterangan
        mengubah \omterm{def:g10:proba:distribution}{peluang}, dan cara keterangan itu diperoleh pun penting}
        (pertanyaan 7--14); \omterm{def:g10:proba:sample}{sampel} berfluktuasi sekitar
        $\frac{1}{\sqrt n}$, dan simulasi yang menengahi
        (pertanyaan 11, 16--18).

\end{enumerate}
\end{problem}
